\element{1.5.1}{
  \begin{question}{modulationAmplitude}\bareme{1}
    Un signal modulé en amplitude s'écrit
    \begin{multicols}{2}
      \begin{reponses}
        \bonne{$s(t)=A(t)\cos(2\pi ft+\phi)$}
        \mauvaise{$s(t)=A\cos(2\pi f(t)t+\phi)$}
        \mauvaise{$s(t)=A\cos(2\pi ft+\phi(t))$}
        \mauvaise{$s(t)=A\cos(2\pi f(t)+\phi)$}
      \end{reponses}
    \end{multicols}
  \end{question}
}
\element{1.5.1}{
  \begin{question}{modulationFréquence}\bareme{1}
    Un signal modulé en fréquence s'écrit
    \begin{multicols}{2}
      \begin{reponses}
        \mauvaise{$s(t)=A(t)\cos(2\pi ft+\phi)$}
        \bonne{$s(t)=A\cos(2\pi f(t)t+\phi)$}
        \mauvaise{$s(t)=A\cos(2\pi ft+\phi(t))$}
        \mauvaise{$s(t)=A\cos(2\pi f(t)+\phi)$}
      \end{reponses}
    \end{multicols}
  \end{question}
}
\element{1.5.1}{
  \begin{question}{modulationPhase}\bareme{1}
    Un signal modulé en phase s'écrit
    \begin{multicols}{2}
      \begin{reponses}
        \mauvaise{$s(t)=A(t)\cos(2\pi ft+\phi)$}
        \mauvaise{$s(t)=A\cos(2\pi f(t)t+\phi)$}
        \bonne{$s(t)=A\cos(2\pi ft+\phi(t))$}
        \mauvaise{$s(t)=A\cos(2\pi f(t)+\phi)$}
      \end{reponses}
    \end{multicols}
  \end{question}
}
\setgroupmode{1.5.1}{cyclic}

\element{1.5.2}{
  \begin{questionmult}{demodulation}\bareme{haut=2,p=0}
    Parmi les opérations ci-dessous, lesquelles sont effectuées pour retrouver le signal modulant avec la démodulation synchrone ?
    \begin{multicols}{2}
      \begin{reponses}
        \bonne{La multiplication par la porteuse}
        \mauvaise{La multiplication par le signal à transmettre}
        \bonne{Un filtrage passe-haut}
        \bonne{Un filtrage passe-bas}
        \mauvaise{Un filtrage passe-bande}
        \mauvaise{La soustraction de la porteuse}
      \end{reponses}
    \end{multicols}
  \end{questionmult}
}

\element{1.5.3}{
  \begin{questionmult}{spectre}\bareme{haut=1,p=0}
    On module en amplitude un signal $s(t)=A\cos(2\pi f t)$ par une porteuse $s_p(t)=A_p\cos(2\pi f_pt)$.
    Parmi les fréquences suivantes, lesquelles sont présentent dans le spectre du signal modulé ?
    \begin{multicols}{3}
      \begin{reponses}
        \bonne{$f_p$}
        \bonne{$f_p-f$}
        \bonne{$f_p+f$}
        \mauvaise{$0$}
        \mauvaise{$2f_p$}
        \mauvaise{$2f$}
        \mauvaise{$2f_p-f$}
        \mauvaise{$2f_p+f$}
        \mauvaise{$f$}
      \end{reponses}
    \end{multicols}
  \end{questionmult}
}

\element{1.5.4}{
  \begin{question}{formuleTrigo}\bareme{1}
    $\cos(a)\cos(b)=$
    \begin{multicols}{2}
      \begin{reponses}
        \bonne{$\frac{\cos(a+b)+\cos(a-b)}{2}$}
        \mauvaise{$\frac{\cos(a+b)-\cos(a-b)}{2}$}
        \mauvaise{$\cos(\frac{a+b}{2})+\cos(\frac{a-b}{2})$}
        \mauvaise{$\cos(\frac{a+b}{2})-\cos(\frac{a-b}{2})$}
        \mauvaise{$\sin(a)\sin(b)-\cos(a)\sin(b)$}
        \mauvaise{$\cos(a)\sin(b)-\sin(a)\sin(b)$}
      \end{reponses}
    \end{multicols}
  \end{question}
}